% v13: the constant-pairing C4 obstruction no longer uses D opposition.
% The new tau=1 argument uses the fifth core and the unique S=T=2 model.

\subsection{四角形成分の同符号 pairing の排除}
\label{cq12:section}

\begin{lemma}[二つの reverse divisors が二である最小 entry の模型]
\label{cq13:reverse-two}
正の二高分母標準形で $S=T=2$ なら，canonical な整数データは
\[
 (r,s,z,S,T)=(4,4,5,2,2),\qquad (a,b)=(9,11)
\]
に限られる．
\end{lemma}

\begin{proof}
$d=S-T-1=-1$ なので，補題~\ref{ml12:dichotomy} より
$\alpha=-u$，$\beta=-v$，$H=r-2\geq1$ と書ける．ここで $u,v$ は正の整数である．
式~\eqref{ml12:E-equation} と式~\eqref{ml12:s-z} は
\[
 Euv=v-u+1,\qquad E=1-\frac4{H(H+1)},\qquad
 s=u\frac{H+2}{H},\quad z=v\frac{H+3}{H+1}
\]
を与える．$H=1$ なら $E=-1$ で，最初の式は $u(1-v)=v+1$ となる．
左辺は非正，右辺は正なので不可能である．
$H\geq2$ では $0<E<1$ である．$u=1$ なら $E=1$ となるため $u\geq2$ であり，
\[
 Eu=1-\frac{u-1}{v}<1
\]
となる．$H\geq3$ なら $E\geq2/3$ なので $Eu\geq4/3$ となって矛盾する．
従って $H=2$，$E=1/3$，$u=2$，$v=3$ である．
表示式から $r=4$，$s=4$，$z=5$，$(a,b)=(sS+1,zT+1)=(9,11)$ を得る．
このデータが標準形の各条件を満たすことは直接確認できる．
\end{proof}

\begin{theorem}[四角形成分に残る pairing]\label{cq12:no-constant-pairing}
\label{cq13:no-constant-pairing}
高分母 support の連結成分が $C_4$ であり，その vertices を
$i\in\mathbb Z/4$，残る第五の label を $*$ とする．
正の局所 charts，marked puncture の整数 clutching，
および全ての node 数の非負性を仮定する．
\[
 \varepsilon_i=\operatorname{sign}(p^i_{i+1,*})
\]
が全て同じであることは不可能である．
従って命題~\ref{cn12:cycle-signs} の二つの型のうち，残る可能性は
$\varepsilon_i=-\varepsilon_{i+1}$ という交互型だけである．
ここで node 数の上限五は，別の仮定ではなく両端の row sums と非負性から従う．
\end{theorem}

\begin{proof}
cycle の向きを逆にすると全 $\varepsilon_i$ の符号が反転するので，
全て $+1$ としてよい．$q_i=q_{i,i+1}>1$ とする．
core $i$ の高分母二 rays 間の entries を
\[
 p^i_{i-1,i+1}=-A_i,\qquad
 p^i_{i+1,i-1}=-B_i,
 \qquad A_iB_i=q_{i-1}q_i-1
\]
と書く．全ての $A_i,B_i$ は正である．

\paragraph{固定向きの低低 entry と reverse divisors．}
共通 puncture $*$ では $\lambda_{i,i+1}=0$ である．
その reverse numerators の絶対値を $c_i$ とすると，補題~\ref{cn12:reverse-bound} より
$1\leq c_i\leq4$ である．
core $i$ の低 ray $i*$ から二本の高 rays へ向かう entries は，順に
\[
 p^i_{*,i-1}=-c_{i-1},\qquad p^i_{*,i+1}=c_i.
\]
もう一つの低 ray $i,i+2$ から高 rays への entries を
\[
 p^i_{i+2,i-1}=\ell_i>0,\qquad
 p^i_{i+2,i+1}=-k_i<0
\]
と書く．固定向きの低低 entries は
\[
 d_i=p^i_{*,i+2}\leq0,\qquad e_i=p^i_{i+2,*}\leq0
\]
であり，ちょうど一方が零，他方が負である．
特に $d_i$ は canonical な非零 entry と同じとは限らない．
低 ray $i*$ の row sum から
\begin{equation}\label{cq12:fixed-row}
 c_i-c_{i-1}+d_i=-1,
 \qquad \sum_i d_i=-4
\end{equation}
を得る．

\paragraph{mixed node による二つの型の伝播．}
$ij=(i,i+1)$ の，第三 label $i+2$ に向かう mixed puncture の両端の outgoing entries は
\[
 -\frac{q_i-1}{k_i},\qquad -A_{i+1}
\]
であり，両者とも負である．従ってその node 数は二以上の整数である．
これを $1+\tau_i$ と書く．共通 puncture の node 数は一であり，もう一つの mixed puncture の node 数も一以上なので
\[
 \tau_i\in\{1,2\}.
\]
全数五を用いて，二つの mixed clutching 式は
\begin{equation}\label{cq12:tau-equations}
 A_{i+1}=\tau_iq_i-\frac{q_i-1}{k_i},\qquad
 B_i=(2-\tau_i)q_i+\frac{q_i-1}{\ell_{i+1}}
\end{equation}
となる．各表示分数は一以上 $q_i-1$ 以下の整数なので，
\[
 \begin{array}{c|cc}
  & A_{i+1} & B_i\\ \hline
 \tau_i=1 & <q_i & >q_i\\
 \tau_i=2 & >q_i & <q_i
 \end{array}
\]
である．一方 $A_iB_i=q_{i-1}q_i-1$ と整数性により，
$A_i<q_{i-1}$ と $B_i>q_i$ は同値であり，逆向きの厳密な不等式も同値である．
従って $\tau_{i-1}=\tau_i$ であり，四つの $\tau_i$ は同じ値となる．

\paragraph{全て $\tau_i=1$ の場合．}
$d_i<0$ の core を考える．このとき $e_i=0$ なので，二高分母の canonical 表示の低 ray $C$ は $i,i+2$，
低 ray $D$ は $i*$ である．$C$ の正の entry は $i,i-1$ に向かうので，
canonical な高 ray $A$ は $i,i-1$，$B$ は $i,i+1$ となる．
canonical 補助変数は
\[
 \alpha=A_i-q_{i-1}<0,\qquad
 \beta=q_i-B_i<0
\]
である．補題~\ref{ml12:dichotomy} より，その非零低低 entry $d_i$ は $-1$ である．
従って全ての $d_i$ は零または $-1$ であり，\eqref{cq12:fixed-row} から四つとも $-1$ となる．
同じ式から $c_i=c$ は一定である．この場合の canonical な低 rays は全て
$C=i+2$，$D=*$ であり，$S=T=c$ となる．
補題~\ref{cc12:proper-negative} の結論 $S\geq2$ により $c\geq2$ である．
以下では分母の下降や $D$ の相殺を使わず，第五の core $*$ の node 条件を用いる．

core $*$ における，ray $*i$ から ray $*j$ への entry を $v_{ij}=p^*_{ij}$ と書く．
これらの ray の分母は全て一なので
\[
 \sum_{j\ne i}v_{ij}=-1,\qquad v_{ij}v_{ji}=0.
\]
低 edge $i*$ の puncture $i+1$ では，core $i$ 側の entry は $c$ である．
従って node 非負性は
\[
 1-c-v_{i,i+1}\geq0,\qquad v_{i,i+1}\leq1-c<0
\]
を与える．添字を一つ戻して pair product を使えば $v_{i,i-1}=0$ であり，row sum から
\[
 v_{i,i+2}=-1-v_{i,i+1}\geq c-2
\]
となる．$c\geq3$ なら $v_{i,i+2}$ と $v_{i+2,i}$ がともに正となり，
その積が零であることに反する．従って $c=2$ である．

補題~\ref{cq13:reverse-two} により，各 core の canonical 高分母は $(a,b)=(9,11)$ である．
ここでは $A=i-1$，$B=i+1$ なので，すべての $i$ に対して
$q_{i-1}=9$，$q_i=11$ を要求する．同じ共有 edge の分母が両方の値を持つことになり矛盾する．

\paragraph{全て $\tau_i=2$ の場合．}
同じ標識の説明により，$d_i<0$ の core では canonical な
$\alpha=A_i-q_{i-1}$，$\beta=q_i-B_i$ がともに正である．
従って $d_i\leq-3$ である．
四つの和が $-4$ なので，ちょうど一つが $-4$，他の三つが零である．
添字を移して $d_0=-4$ とすれば，\eqref{cq12:fixed-row} と $1\leq c_i\leq4$ から
\[
 (c_0,c_1,c_2,c_3)=(4,3,2,1)
\]
となる．
core $3$ では $d_3=0$ なので，canonical な低 ray $C$ は $3*$ であり，
その正の entry は高 ray $3,0$ に向かう．従って canonical な高 ray $A$ は $3,0$ で，
標準形の整数 $r$ は $c_3=1$ である．
ところが $r=1$ の標準形では
\[
 A=a+\frac{a-1}{r}-s=a+(a-1)-s\geq a,
 \qquad s\mid a-1,\quad s>0.
\]
これは $\tau_3=2$ から得る $B_3<q_3$ に反する．
ここでは canonical な $A$ が $B_3$，$a$ が $q_3$ に対応する．
両場合が排除され，結論を得る．
\end{proof}

交互型の四角形成分については，本定理は不存在を結論しない．
同符号型の排除には node 非負性と各 branch の全 node 数五を使っており，
node 条件を外した整数 clutching だけの障害へ読み替えない．
一方で，この改訂された証明では対向端の $D$ の相殺を一切用いていない．
